\documentclass[11pt, a4paper, oneside]{article}
\usepackage[ngerman]{babel}
\usepackage{worksheet}

\title{Webentwicklung: HTML}
\author{L. Bung}
\subject{SAE}
\class{E3FI}

\begin{document}
	\maketitle
	
	\begin{hint}{Was ist HTML?}
		HTML (Hypertext Markup Language) ist die standardisierte Auszeichnungssprache zur Strukturierung digitaler Dokumente im World Wide Web.
		HTML beschreibt \textbf{nicht} das Aussehen (Design) einer Seite, sondern ausschließlich deren logischen Aufbau und die Bedeutung der Inhalte (Semantik).
	
		Das Grundgerüst einer HTML-Seite ist folgendermaßen aufgebaut:
		
		\begin{lstlisting}[language=HTML, numbers=none, frame=none]
<!DOCTYPE html>
<html lang="de">
	<head>
		<meta charset="UTF-8">
		<title>Meine erste Webseite</title>
	</head>
	<body>
		<!-- Hier steht der sichtbare Inhalt der Webseite -->
	</body>
</html>
		\end{lstlisting}
	\end{hint}
	
	\singletask{Grundlegende Textstrukturierung}
	
	\begin{enumerate}[label=\alph*)]
		\item Integrieren Sie in den \texttt{<body>} des oben stehenden Grundgerüsts einen kurzen Steckbrief über Ihren Ausbildungsbetrieb (oder ein fiktives Unternehmen).
		Nutzen Sie dabei zwingend die folgenden HTML-Tags, um den Text logisch zu strukturieren:
		
		\begin{itemize}
			\item Hauptüberschrift (\texttt{<h1>})
			\item Unterüberschriften (\texttt{<h2>}, \texttt{<h3>} usw. bis \texttt{<h6>})
			\item Absätze (\texttt{<p>})
			\item eine ungeordnete Liste (\texttt{<ul>})
			\item eine geordnete Liste (\texttt{<ol>})
		\end{itemize}
		
		\item Finden Sie heraus, was im \texttt{<head>}-Teil eines HTML-Dokuments steht und wofür diese Informationen benötigt werden.
	\end{enumerate}
	
	\lines[3cm]
	
	
	\partnertask{HTML-Attribute, Verlinkungen und Medien}
	
	\begin{enumerate}[label=\alph*)]
		\item Recherchieren Sie, was man bei HTML unter Attributen versteht.
		Nennen Sie Beispiele für häufig verwendete Attribute und erklären Sie, wofür man sie benötigt.
		\item Erweitern Sie Ihren Steckbrief aus Aufgabe 1 um einen Link zur Firmenwebsite (\texttt{<a>}) sowie ein passendes Firmenlogo (\texttt{<img>}).
		\item Erklären Sie stichpunktartig die Bedeutung und Notwendigkeit der Attribute \texttt{href}, \texttt{src} und \texttt{alt} in Bezug auf die letzte Teilaufgabe.
	\end{enumerate}
	\lines[5cm]
	
	
	\singletask{Semantisches Markup}
	
	Früher wurden Webseiten fast ausschließlich mit generischen Containern (\texttt{<div>}) strukturiert.
	
	\begin{enumerate}[label=\alph*)]
		\item Diskutieren Sie, warum dies problematisch sein könnte.
		\item Bereinigen Sie die nachfolgende HTML-Struktur.
		Finden Sie dazu passende Tags, die die generischen Container ersetzen.
	\end{enumerate}
	
	\begin{lstlisting}[language=HTML, numbers=left, frame=single]
<div id="kopfbereich"> ... </div>
<div id="hauptnavigation"> ... </div>
<div id="inhaltsbereich">
	<div class="blog-eintrag"> ... </div>
	<div class="blog-eintrag"> ... </div>
</div>
<div id="seitenleiste"> ... </div>
<div id="fussbereich"> ... </div>
	\end{lstlisting}
	
	\pagebreak
	
	\grouptask{Trennung von Struktur und Design}
	
	In älteren HTML-Versionen gab es Tags wie \texttt{<font>} und \texttt{<center>} oder Attribute wie \texttt{bgcolor}, um Text rot zu färben oder zu zentrieren.
	In modernem HTML5 sind diese veraltet (deprecated) und sollen nicht mehr verwendet werden.
	Das Design wird heute vollständig in CSS (Cascading Style Sheets) ausgelagert -- mehr dazu in einer der kommenden Stunden.
	
	Diskutieren Sie in einer 3er- bis 4er-Gruppe:
	\begin{itemize}
		\item Welche konkreten Nachteile entstehen bei der Softwareentwicklung und Wartung, wenn Design und Struktur im HTML-Code eng miteinander vermischt werden?
		\item Fallen Ihnen Beispiele aus der Programmierung ein, bei denen Sie Ähnlichkeiten zu diesem Problem sehen können?
		\item In welchen Fällen könnte sogenanntes Inline-Styling (Design direkt im HTML-Code) heute noch pragmatisch oder sogar notwendig sein?
	\end{itemize}
	Halten Sie die wichtigsten Argumente Ihrer Gruppe stichpunktartig fest.
	
	\lines[8cm]
	
	
	\bonustask{Erweiterung des Steckbriefs}
	
	Erweitern Sie den Steckbrief Ihres Ausbildungsbetriebs um weitere HTML-Seiten.
	Finden Sie heraus, wie man zwischen zwei lokalen HTML-Seiten verlinken kann.
	Nutzen Sie dies anschließend, um die verschiedenen Dateien zu verknüpfen -- achten Sie darauf, dass man von jeder Seite aus durch Links auf alle anderen Seiten kommen kann.
\end{document}